% =====================================================================
%  Du cadrage à l'interface : TP1, TP2, TP3, version détaillée
%  Module : Développement et déploiement d'applications intelligentes
%  A. Larhlimi
%  Compilation : pdflatex tp1-tp2-tp3-detaille.tex   (deux fois)
% =====================================================================
\documentclass[11pt,a4paper]{article}

\usepackage[T1]{fontenc}
\usepackage[utf8]{inputenc}
\usepackage[french]{babel}
\IfFileExists{lmodern.sty}{\usepackage{lmodern}}{}
\usepackage[margin=2cm]{geometry}
\usepackage{booktabs}
\usepackage{longtable}
\usepackage{array}
\usepackage{xcolor}
\usepackage{enumitem}
\usepackage{listings}
\usepackage{fancyhdr}
\usepackage{etoolbox}
\usepackage{titlesec}
\usepackage{tcolorbox}
\IfFileExists{xurl.sty}{\usepackage{xurl}}{}
\IfFileExists{hyphenat.sty}{\usepackage[htt]{hyphenat}}{}
\usepackage[hidelinks]{hyperref}

\definecolor{navy}{HTML}{1B2A41}
\definecolor{orangeacc}{HTML}{E07A2F}
\definecolor{teal}{HTML}{2A9D8F}
\definecolor{grisclair}{HTML}{F0EDE5}
\definecolor{grismoyen}{HTML}{5B6572}

\newcommand{\co}[1]{\texttt{\detokenize{#1}}}
\newcommand{\brk}{\discretionary{}{}{}}

% Encadrés : objectif, contrôle, piège
\newtcolorbox{objectif}[1][]{colback=teal!6, colframe=teal, boxrule=1pt, arc=2pt,
  left=6pt, right=6pt, top=4pt, bottom=4pt, title={\bfseries Objectif}, fonttitle=\color{white}, #1}
\newtcolorbox{controle}[2][]{colback=teal!6, colframe=teal, boxrule=1pt, arc=2pt,
  left=6pt, right=6pt, top=4pt, bottom=4pt, title={\bfseries Contrôle #2}, fonttitle=\color{white}, #1}
\newtcolorbox{piege}[2][]{colback=orangeacc!8, colframe=orangeacc, boxrule=1pt, arc=2pt,
  left=6pt, right=6pt, top=4pt, bottom=4pt, title={\bfseries #2}, fonttitle=\color{white}, #1}

\lstdefinestyle{tp}{
  basicstyle=\ttfamily\footnotesize,
  backgroundcolor=\color{grisclair},
  frame=leftline, framerule=1.2pt, rulecolor=\color{teal},
  xleftmargin=8pt, framexleftmargin=8pt,
  breaklines=true, showstringspaces=false,
  commentstyle=\color{grismoyen}\itshape,
  keywordstyle=\color{navy}\bfseries,
  extendedchars=true,
  literate={é}{{\'e}}1 {è}{{\`e}}1 {ê}{{\^e}}1 {à}{{\`a}}1 {â}{{\^a}}1
           {î}{{\^i}}1 {ï}{{\"i}}1 {ô}{{\^o}}1 {û}{{\^u}}1 {ù}{{\`u}}1
           {ç}{{\c c}}1 {É}{{\'E}}1 {À}{{\`A}}1 {Ô}{{\^O}}1 {È}{{\`E}}1
           {«}{{\guillemotleft}}1 {»}{{\guillemotright}}1 {·}{{\textperiodcentered}}1
           {├}{{\textbar}}1 {└}{{\textbar}}1 {│}{{\textbar}}1 {─}{{-}}1 {’}{{'}}1
}
\lstset{style=tp}

\pagestyle{fancy}
\fancyhf{}
\fancyhead[L]{\small TP1, TP2, TP3 \textperiodcentered{} Du cadrage à l'interface}
\fancyhead[R]{\small A. Larhlimi}
\fancyfoot[C]{\small\thepage}
\renewcommand{\headrulewidth}{0.4pt}

\titleformat{\section}{\normalfont\Large\bfseries\color{navy}}{\thesection}{0.6em}{}
\titleformat{\subsection}{\normalfont\large\bfseries}{\thesubsection}{0.6em}{}

\setlist[itemize]{topsep=3pt,itemsep=2pt,leftmargin=*}
\setlist[enumerate]{topsep=3pt,itemsep=2pt,leftmargin=*}
\renewcommand{\arraystretch}{1.25}
\AtBeginEnvironment{longtable}{\small}
\setlength{\tabcolsep}{4pt}
\setlength{\emergencystretch}{3em}
\sloppy

\title{\vspace{-1.4cm}\textbf{Du cadrage à l'interface}\\[4pt]
\large TP1, TP2, TP3 : construire une application de sentiment, de A à Z}
\author{A. Larhlimi}
\date{Développement et déploiement d'applications intelligentes\\
Trois séances de 2 h \textperiodcentered{} travail en binôme \textperiodcentered{} \today}

\begin{document}
\maketitle
\thispagestyle{fancy}

\noindent
Ce document enchaîne les trois premières séances sur un seul et même dépôt. Il donne, pour chaque
étape, son objectif, les commandes avec leur explication, le résultat attendu et les erreurs
possibles. À la fin, votre binôme dispose d'une application complète : un modèle entraîné et évalué,
et une interface qu'une personne non technique peut utiliser.

\tableofcontents
\bigskip

% =====================================================================
\section{Vue d'ensemble}
% =====================================================================

\subsection{Le besoin}

Une équipe relation client reçoit des milliers d'avis par jour et les traite dans l'ordre d'arrivée.
On veut faire remonter les \textbf{avis négatifs} en tête de file, pour répondre d'abord aux clients
mécontents.

Le problème est volontairement simple côté modèle. Classer un avis en positif ou négatif se fait en
vingt lignes de code, sans GPU, avec un classifieur linéaire. La difficulté est ailleurs, et c'est
le sujet du module : rendre ce modèle reproductible, testé, surveillé, utilisable par quelqu'un
d'autre, et redéployable après un réentraînement. Autrement dit, \textbf{l'écart entre un notebook
et un produit}.

\subsection{Les trois séances}

\begin{center}
\begin{tabular}{@{}p{1.5cm}p{8.4cm}p{2.4cm}p{2.2cm}@{}}
\toprule
\textbf{Séance} & \textbf{Ce que le binôme produit} & \textbf{Contrôles} & \textbf{Tests} \\
\midrule
TP1 & Dépôt structuré, environnement reproductible, trois fiches de cadrage & T0 à T7 & 38 \\
TP2 & Le modèle : entraînement, évaluation, métriques versionnées & C1 à C5 & 42 \\
TP3 & L'interface de démonstration Streamlit & S1 à S4 & 45 \\
\bottomrule
\end{tabular}
\end{center}

Un seul dépôt, \co{sentiment-app}, du premier jour au dernier. Chaque kit se pose sur le précédent
sans rien écraser, et les tests s'accumulent au lieu de se remplacer.

\subsection{Les compétences visées}

\begin{center}
\begin{tabular}{@{}p{9.0cm}p{6.4cm}@{}}
\toprule
\textbf{Compétence} & \textbf{Vérifiée par} \\
\midrule
Choisir une approche technique et la justifier & 12 tests de cadrage, lecture des fiches \\
Produire un modèle reproductible et l'évaluer correctement & 4 tests de modèle, \co{metrics.json} \\
Exposer le modèle à un utilisateur non technique & 3 tests d'interface \\
Travailler à deux sur un dépôt, avec tests et intégration continue & 23 tests d'environnement, 3 d'historique \\
\bottomrule
\end{tabular}
\end{center}

\subsection{Comment lire ce document}

Dix-sept contrôles numérotés jalonnent les trois séances : T0 à T7 pour le TP1, C1 à C5 pour le TP2,
S1 à S4 pour le TP3. À chacun, le résultat attendu est écrit noir sur blanc.

\textbf{Ne passez à l'étape suivante que lorsque le contrôle précédent a donné ce résultat.} Deux
d'entre eux doivent être rouges : c'est précisé à l'endroit voulu. Un binôme qui enchaîne sans
vérifier arrive à la fin avec trois problèmes mélangés et ne sait plus lequel vient d'où.

% =====================================================================
\section{Prérequis et réglages}
% =====================================================================

\begin{objectif}
Arriver avec un poste prêt. Les quatre réglages de cette section suppriment à eux seuls trois des
blocages les plus fréquents des séances.
\end{objectif}

\subsection{Ce qu'il faut installer}

\begin{center}
\begin{tabular}{@{}p{3.4cm}p{3.0cm}p{9.0cm}@{}}
\toprule
\textbf{Outil} & \textbf{Version} & \textbf{Remarque} \\
\midrule
Python & 3.11.x (3.10 toléré) & Windows : cocher « Add python.exe to PATH ». \textbf{Évitez 3.12 et 3.13} : les versions figées du cours n'ont pas de paquet précompilé pour elles \\
Git & dernière stable & L'installeur Windows fournit aussi Git Bash, qui exécute les scripts \co{.sh} \\
VS Code & dernière & Facultatif, mais le kit contient sa configuration \\
Compte GitHub & & Avec l'adresse de l'école, de préférence \\
\bottomrule
\end{tabular}
\end{center}

\subsection{Vérifier, puis régler}

\begin{lstlisting}[language=bash]
python --version                    # 3.11.x attendu ; sinon py -3.11 --version
git --version                       # si « commande inconnue », rouvrez le terminal
\end{lstlisting}

\begin{lstlisting}[language=bash]
git config --global user.name "Prénom Nom"
git config --global user.email "prenom.nom@ecole.ma"
git config --global init.defaultBranch main
Set-ExecutionPolicy -Scope CurrentUser RemoteSigned
\end{lstlisting}

Ces quatre lignes méritent une explication, car elles reviennent dans trois cas problèmes.

\textbf{\co{user.name}} est inscrit comme auteur dans chacun de vos commits. C'est ce nom, et non
l'adresse, que compte le contrôle T5. Sans identité configurée, Git refuse de créer le commit avec
\co{Please tell me who you are}.

\textbf{\co{user.email}} sert à GitHub pour rattacher le commit à votre compte. Si elle ne correspond
à aucune adresse enregistrée dans votre profil, vos commits apparaissent sans votre photo ni votre
lien, même s'ils sont bien de vous. Git ne vérifie jamais cette adresse : elle déclare qui vous dites
être, pas qui vous êtes. L'authentification, elle, passe par le jeton d'accès personnel.

\textbf{\co{--global}} écrit dans votre fichier personnel, qui vaut pour tous les dépôts de ce compte
utilisateur. La portée \co{--local}, elle, n'écrit que dans le \co{.git/config} du dépôt courant et
l'emporte sur le global ; elle exige donc d'être \emph{dans} un dépôt, sans quoi Git répond
\co{--local can only be used inside a git repository}.

\textbf{\co{Set-ExecutionPolicy}} autorise l'exécution des scripts PowerShell pour votre seul compte.
Sans elle, \co{Activate.ps1} est refusé et le venv ne s'active pas.

\subsection{Répartir les rôles}

\begin{center}
\begin{tabular}{@{}p{2.4cm}p{6.4cm}p{6.6cm}@{}}
\toprule
\textbf{Rôle} & \textbf{Ce qu'il fait seul} & \textbf{Ce qu'il fait comme l'autre} \\
\midrule
Membre A & Crée le dépôt GitHub, pousse le kit & venv, installation, son commit, vérifications \\
Membre B & Clone le dépôt une fois créé & venv, installation, son commit, vérifications \\
\bottomrule
\end{tabular}
\end{center}

Décidez maintenant : les étapes 4 et 5 du TP1 en dépendent.

% =====================================================================
\section{Les ressources partagées}
% =====================================================================

\subsection{Trois kits qui s'empilent}

\begin{center}
\begin{longtable}{@{}p{4.4cm}p{6.6cm}p{4.4cm}@{}}
\toprule
\textbf{Archive} & \textbf{Contenu} & \textbf{Se pose sur} \\
\midrule
\endhead
\co{sentiment-app.zip} & Le dépôt de départ complet : structure, scripts, 38 tests, CI, six documents & rien, c'est la base \\
\co{tp2-semaine2-kit.zip} & \co{src/} (quatre scripts), \co{tests/test_model.py}, notebook Colab & le dépôt du TP1 \\
\co{tp3-streamlit-kit.zip} & \co{ui/} (interface et helpers), \co{tests/test_ui.py}, notebook Colab, documents du compte rendu & le dépôt après le TP2 \\
\bottomrule
\end{longtable}
\end{center}

\subsection{Installer un kit}

Décompressez l'archive \textbf{à côté} de votre dépôt, pas dedans. Vous devez obtenir deux dossiers
voisins :

\begin{lstlisting}[language=bash]
IA-DEV/
├── sentiment-app/      <- votre dépôt
└── tp2-semaine2/       <- le kit que vous venez de décompresser
\end{lstlisting}

\begin{lstlisting}[language=bash]
cd tp2-semaine2
python appliquer_tp2.py ../sentiment-app --simulation   # voir sans rien écrire
python appliquer_tp2.py ../sentiment-app                # copier pour de bon
\end{lstlisting}

Le script vérifie d'abord que la cible ressemble bien au dépôt attendu, puis copie les fichiers un
par un en annonçant pour chacun s'il l'ajoute ou le remplace. Celui du TP2 ajoute en outre une ligne
au \co{.gitignore}, celui du TP3 refuse de s'installer si \co{src/predict.py} est absent, c'est-à-dire
si le TP2 n'a pas été fait.

\subsection{L'entrepôt complet}

\begin{center}
\begin{longtable}{@{}p{3.2cm}p{5.6cm}p{3.0cm}p{3.4cm}@{}}
\toprule
\textbf{Dossier} & \textbf{Fichier} & \textbf{Pour qui} & \textbf{Quand} \\
\midrule
\endhead
 & \co{tp1-tp2-tp3.md} et ce PDF & Tous & La démarche complète \\
 & \co{TP3-colab-streamlit.ipynb} & Étudiants & L'interface sur Colab \\
\co{kits/} & les trois archives & Étudiants & Une par séance \\
\co{compte-rendu/} & énoncé et modèle à remplir & Étudiants & Après le TP3 \\
\co{diapositives/} & quatre PDF & Projection & Vue d'ensemble, TP1 détaillé, TP1 démarche, TP2 \\
\co{enseignant/} & corrigés, QCM, guide & \textbf{Enseignant seulement} & Ne pas distribuer \\
\bottomrule
\end{longtable}
\end{center}

Calendrier de diffusion conseillé : les prérequis une semaine avant, le kit de la séance deux jours
avant, les diapositives le jour même, l'énoncé du compte rendu après le TP3.

% =====================================================================
\section{Architecture du dépôt}
% =====================================================================

\begin{objectif}
Savoir à quoi sert chaque fichier avant de l'utiliser. Un dépôt dont on ne comprend pas la structure
se dégrade en trois semaines.
\end{objectif}

\subsection{Les sept dossiers}

\begin{center}
\begin{tabular}{@{}p{2.8cm}p{5.6cm}p{7.0cm}@{}}
\toprule
\textbf{Dossier} & \textbf{Contient} & \textbf{Se remplit} \\
\midrule
\co{data/} & Jeux de données & TP2, versionné par DVC en semaine 7 \\
\co{models/} & Modèles entraînés & TP2. Jamais dans Git : lourd, reconstructible \\
\co{app/} & API FastAPI & Semaine suivante \\
\co{ui/} & Interface Streamlit & TP3 \\
\co{tests/} & Tests pytest & Dès le TP1 \\
\co{scripts/} & Installation et diagnostic & Fourni \\
\co{docs/cadrage/} & Fiches de cadrage & Partie A du TP1 \\
\bottomrule
\end{tabular}
\end{center}

Les dossiers encore vides contiennent un fichier \co{.gitkeep}, vide lui aussi. Raison : \textbf{Git
ne versionne pas les dossiers, seulement les fichiers}. Sans ce témoin, un dossier vide disparaîtrait
au clone, et la structure ne serait pas la même chez tout le monde.

\subsection{requirements.txt}

\begin{lstlisting}[language=bash]
pandas==2.2.*
scikit-learn==1.5.*
joblib==1.4.*
pyarrow==21.0.*
fastapi==0.115.*
streamlit==1.39.*
mlflow==2.17.*
\end{lstlisting}

La syntaxe \co{==2.2.*} accepte 2.2.0 ou 2.2.3, jamais 2.3 : les correctifs passent, les changements
de comportement non. Le test \co{test_requirements_pins_versions} vérifie que chaque ligne utile
contient \co{==}.

Pourquoi installer tout le semestre dès la première séance ? Pour ne plus jamais y revenir. Une
installation de dix minutes aujourd'hui vaut mieux que dix interruptions de deux minutes réparties
sur onze semaines.

\subsection{.gitignore}

\begin{lstlisting}[language=bash]
.venv/                     # l'environnement virtuel : des centaines de Mo
__pycache__/               # le bytecode, régénéré à chaque exécution
.env                       # les secrets : clés d'API, mots de passe
models/*                   # les modèles entraînés
!models/.gitkeep           # exception : le fichier témoin
!models/metrics.json       # exception : la mesure, ajoutée par le kit TP2
data/*                     # les jeux de données
!data/sample_reviews.csv   # exception : le petit jeu d'exemple
.vscode/*
!.vscode/settings.json     # exception : la configuration partagée
\end{lstlisting}

Le \co{!} inverse une règle. Cette mécanique d'exception est le cœur du dépôt : elle sépare ce qui est
\textbf{lourd et reconstructible} (exclu) de ce qui est \textbf{léger et partagé} (inclus). La
section~\ref{sec:mlops} y revient, car c'est elle qui porte tout le raisonnement MLOps.

\subsection{Les autres fichiers}

\begin{center}
\begin{longtable}{@{}p{4.0cm}p{11.6cm}@{}}
\toprule
\textbf{Fichier} & \textbf{Rôle} \\
\midrule
\endhead
\co{.env.example} & Modèle des variables secrètes, versionné car vide de vrais secrets. Chacun le copie en \co{.env}, lui-même ignoré \\
\co{pytest.ini} & \co{testpaths = tests} évite de parcourir tout le dépôt ; \co{addopts = -ra} résume échecs et tests ignorés avec leur raison \\
\co{.editorconfig} & UTF-8, fins de ligne Unix, indentation : évite les conflits illisibles entre Windows et macOS \\
\co{.vscode/settings.json} & Sélectionne l'interpréteur \co{.venv} et active pytest : supprime la cause la plus fréquente du \co{ModuleNotFoundError} \\
\co{Makefile} & Raccourcis \co{make setup}, \co{make check}, \co{make test} (Linux, macOS, Git Bash) \\
\co{README.md} & Porte d'entrée : installation en trois commandes, structure, tableau du binôme \\
\co{MARCHE-A-SUIVRE.md} & Le pas à pas de la séance \\
\co{TESTS.md}, \co{AIDE-GIT.md}, \co{DEPANNAGE.md} & Tests en détail, commandes Git, messages d'erreur \\
\bottomrule
\end{longtable}
\end{center}

% =====================================================================
\section{TP1 : cadrer et préparer l'environnement}
% =====================================================================

\begin{objectif}
Décider quelle approche technique employer pour trois besoins métier, et monter un dépôt que
l'enseignant peut cloner sur une machine vierge et faire tourner en trois commandes.
\end{objectif}

\subsection{Étape 1. Décompresser le kit (5 min, chacun)}

\begin{lstlisting}[language=bash]
unzip sentiment-app.zip        # Windows : clic droit sur le zip, Extraire tout
cd sentiment-app               # toutes les commandes du TP partent d'ici
ls
ls -Force | findstr gitignore  # Linux, macOS : ls -a | grep gitignore
\end{lstlisting}

\begin{controle}{T0 : le kit est complet}
Sept dossiers, six documents, \co{requirements.txt}, \co{pytest.ini}, et \co{.gitignore}, qui ne
s'affiche pas avec un \co{ls} ordinaire parce qu'il commence par un point.
\end{controle}

Pourquoi ce contrôle existe : \co{.gitignore} est le seul rempart contre l'erreur irréversible du
contrôle T3. Son absence se repère maintenant, pas dans vingt minutes.

Si une commande dit plus tard qu'un fichier est introuvable, le premier réflexe est de taper \co{cd}
sans argument pour voir où vous êtes. L'archive crée un dossier \co{sentiment-app} : un chemin du
type \co{sentiment-app\sentiment-app} est normal, descendez jusqu'à celui qui contient
\co{README.md}.

\subsection{Étape 2. L'environnement virtuel (10 min, chacun)}

Le Python du système est unique et partagé par tous vos projets. Les paquets du cours sont figés à
des versions précises : les installer globalement casserait vos autres projets, et l'inverse aussi.
Le venv est une copie légère de Python, propre au projet, avec son propre \co{pip} et son propre
\co{site-packages}. Comme il se reconstruit en une commande, il ne part jamais dans Git.

\begin{lstlisting}[language=bash]
py -3.11 -m venv .venv
.\.venv\Scripts\Activate.ps1          # Linux, macOS : source .venv/bin/activate
python -m pip install pandas scikit-learn pytest
python -m pytest -q
\end{lstlisting}

\begin{center}
\begin{tabular}{@{}p{5.6cm}p{10.0cm}@{}}
\toprule
\textbf{Commande} & \textbf{Ce qu'elle fait} \\
\midrule
\co{python -m venv .venv} & Exécute le module \co{venv} livré avec Python ; crée \co{.venv/} \\
\co{Activate.ps1} & Modifie le PATH du \textbf{terminal courant} : l'invite prend le préfixe \co{(.venv)} \\
\co{python -m pip install} & Passe par l'interpréteur courant, donc ne dépend pas du PATH \\
\bottomrule
\end{tabular}
\end{center}

Trois paquets seulement ici, pour atteindre le premier contrôle en une minute. L'installation
complète vient à l'étape 5.

\begin{piege}{T1 doit être ROUGE}
Résultat attendu : \textbf{3 failed, 31 passed, 4 skipped}. Les trois échecs sont les fiches de
cadrage vides, les quatre tests ignorés attendent un dépôt Git qui n'existe pas encore, et les
31 verts confirment que la structure, les paquets et le \co{.gitignore} sont en place.\\[4pt]
Si vous obtenez \co{38 passed} ici, vous testez un autre dossier. Un test rouge avant que le travail
soit fait est un test utile ; un test vert dès le départ ne vérifie rien.
\end{piege}

Un test mérite d'être regardé : \co{test_virtualenv_active} compare \co{sys.prefix} et
\co{sys.base_prefix}. Sur le Python du système, les deux sont égaux ; dans un venv, non. Il existe
pour que l'oubli d'activation se voie tout de suite, au lieu de produire une erreur obscure trois
commandes plus loin.

\subsection{Étape 3. Les trois fiches de cadrage (50 min, à deux)}

\subsubsection{La grille et l'arbre}

\begin{center}
\begin{tabular}{@{}cp{4.0cm}p{10.4cm}@{}}
\toprule
 & \textbf{Critère} & \textbf{La question à se poser} \\
\midrule
(a) & Données étiquetées & A-t-on des exemples annotés, ou peut-on en produire à faible coût ? \\
(b) & Vérifiabilité & La réponse doit-elle être tracée jusqu'à une source ou une règle ? \\
(c) & Coût par requête & Combien d'appels par jour, et que coûte chacun ? \\
(d) & Conséquence d'une erreur & Que se passe-t-il concrètement quand le système se trompe ? \\
\bottomrule
\end{tabular}
\end{center}

Parcourez les questions dans l'ordre ; la première réponse « oui » donne l'approche.

\begin{enumerate}
  \item Logique connue, stable, devant rester explicable ? $\rightarrow$ \textbf{règles métier}
  \item Données étiquetées pour une tâche répétitive ? $\rightarrow$ \textbf{machine learning}
  \item Réponse à tirer de documents existants, et à citer ? $\rightarrow$ \textbf{RAG}
  \item Sinon, et seulement sinon $\rightarrow$ \textbf{modèle génératif seul}
\end{enumerate}

On ne saute jamais directement à la dernière ligne : c'est l'option la plus coûteuse et la moins
vérifiable.

\subsubsection{Les trois cas}

\begin{center}
\begin{tabular}{@{}p{3.6cm}p{4.6cm}p{7.0cm}@{}}
\toprule
\textbf{Fiche} & \textbf{Le cas} & \textbf{Réponse attendue} \\
\midrule
\co{cas1-pharmacie.md} & Ordonnances incomplètes & Règles métier : (a) pas d'étiquettes, (b) explicabilité exigée, (d) erreur grave \\
\co{cas2-reglement.md} & Règlement intérieur & RAG : (b) la réponse cite l'article, (a) pas de questions étiquetées, (c) volume faible \\
\co{cas3-avis-negatifs.md} & Avis clients négatifs & Machine learning : (a) étiquettes gratuites via la note, (c) milliers d'avis, (d) erreur peu grave \\
\bottomrule
\end{tabular}
\end{center}

Douze minutes par cas. Un membre tient le clavier, l'autre parcourt l'arbre à voix haute et conteste ;
inversez à chaque fiche.

\textbf{Deux approches peuvent se défendre pour un même cas : c'est la justification qui est notée,
pas la case cochée.} Le machine learning reste acceptable pour la pharmacie si vous précisez qu'il
faudrait d'abord constituer un jeu étiqueté et conserver la validation du pharmacien.

\subsubsection{Anatomie d'une fiche}

\begin{lstlisting}[language=bash]
## Besoin en une phrase (obligatoire)

> Remplir chaque rubrique en une à trois phrases.     <- consigne, à laisser en place

Signaler au pharmacien toute ordonnance à laquelle il manque un champ obligatoire.

## Approche retenue (obligatoire)

- [x] Règles métier          <- une seule croix dans tout le fichier
- [ ] Machine learning
- [ ] RAG
- [ ] Modèle génératif seul

## Justification (obligatoire)

(a) Aucune ordonnance étiquetée n'existe au départ. (b) Le pharmacien doit pouvoir
expliquer le refus au patient. (d) Une ordonnance validée à tort a une conséquence grave.
\end{lstlisting}

Les autres rubriques obligatoires : utilisateur final, données nécessaires et leur origine, métrique
de succès et seuil, conséquence d'une erreur et validation humaine, risques éthiques. La neuvième,
facultative, est l'approche écartée et pourquoi.

Deux réflexes attendus : la métrique est chiffrée sur un jeu défini (« rappel d'au moins 0,90 sur la
classe négative », pas « de bons résultats »), et dès que l'erreur a une conséquence réelle, vous
écrivez qui valide et à quel moment.

\begin{controle}{T2 : les fiches passent}
\co{python -m pytest tests/test_cadrage.py -q} doit afficher \textbf{12 passed}. Les messages
d'échec nomment la fiche et la rubrique fautives.
\end{controle}

\begin{piege}{L'erreur qui coûte le plus de temps}
Écrire « il y a des données étiquetées » \textbf{sans} la lettre \co{(a)} fait échouer le test. Les
lettres doivent apparaître littéralement dans la justification, et il en faut au moins deux.
\end{piege}

\subsection{Étape 4. Créer le dépôt et pousser (25 min, membre A)}

\subsubsection{Sur github.com}

\begin{enumerate}
  \item \textbf{New repository} : nom \co{sentiment-app}, visibilité \textbf{Private}.
  \item \textbf{Ne cochez rien} dans « Initialize this repository with » : ni README, ni \co{.gitignore}, ni licence. Le kit les fournit, et un dépôt pré-rempli provoquerait un conflit au premier push.
  \item \textbf{Settings}, section Access, \textbf{Collaborators}, bouton \textbf{Add people} : cherchez par identifiant GitHub, ajoutez le membre B en \textbf{Write} et l'enseignant en \textbf{Read}.
  \item Copiez l'URL HTTPS depuis le bouton vert \textbf{Code}.
\end{enumerate}

Le rôle Write autorise le push, le rôle Read ne l'autorise pas, mais les deux permettent de cloner et
de consulter les exécutions de la CI. L'enseignant peut donc noter sans jamais pouvoir modifier votre
travail. \textbf{L'invitation doit être acceptée} : tant que le membre B ne clique pas sur le lien
reçu par courriel, son clone et son push échouent avec \co{403 Forbidden}.

\subsubsection{Sur votre poste}

\begin{lstlisting}[language=bash]
git init          # crée le dépôt local, c'est-à-dire le dossier caché .git/
git add -A        # met tout en attente, sauf ce qu'exclut .gitignore
git status        # CONTRÔLE T3 : lire la sortie AVANT de valider
\end{lstlisting}

\begin{piege}{T3 : le seul contrôle irréversible du TP}
La sortie ne doit contenir \textbf{aucune ligne commençant par \co{.venv/}}. Si elle en contient,
arrêtez-vous : le \co{.gitignore} manque ou a été renommé.\\[4pt]
Une fois le commit poussé, le dépôt pèse 300 Mo, GitHub refuse les fichiers de plus de 100 Mo, et le
test \co{test_venv_never_committed} reste rouge même après un \co{git rm -r --cached .venv}, parce
qu'il lit \textbf{tout l'historique}. La seule correction est alors de refaire le dépôt.
\end{piege}

Profitez-en pour repérer tout fichier parasite : une archive téléchargée, un PDF, une capture d'écran.
Un \co{git rm --cached "le fichier"} et une ligne \co{*.zip} dans le \co{.gitignore} règlent le cas.

\begin{lstlisting}[language=bash]
git commit -m "Semaine 1 : structure du depot et kit TP1"
git branch -M main                 # GitHub attend main ; -M force le renommage
git remote add origin https://github.com/votre-compte/sentiment-app.git
git remote -v                      # vérifiez : PAS de chevrons dans l'URL
git push -u origin main            # -u mémorise le lien : ensuite git push suffit
\end{lstlisting}

Un raccourci fourni enchaîne ces commandes et refuse de pousser si \co{.venv/} est en attente :
\co{bash scripts/init_depot.sh <url>}.

\subsubsection{L'authentification}

Au premier push, Git demande un identifiant et un mot de passe. \textbf{Le mot de passe de votre
compte GitHub est refusé} : il faut un jeton d'accès personnel. GitHub, photo de profil, Settings,
\textbf{Developer settings}, Personal access tokens, \textbf{Tokens (classic)}, Generate new token
(classic), portée \co{repo}. Copiez-le immédiatement, il ne s'affiche qu'une fois, et collez-le à la
place du mot de passe.

Plus simple, si vous préférez éviter les jetons :

\begin{lstlisting}[language=bash]
winget install --id GitHub.cli
gh auth login          # GitHub.com, HTTPS, authentification par navigateur
gh auth setup-git
\end{lstlisting}

\begin{piege}{Deux erreurs qui se ressemblent}
\co{Password authentication is not supported} est un problème d'\textbf{authentification} : GitHub ne
sait pas qui vous êtes. Réponse : le jeton.\\[3pt]
\co{403 Forbidden} est un problème d'\textbf{autorisation} : GitHub le sait, mais vous n'avez pas le
droit d'écrire. Réponse : l'invitation de collaborateur, envoyée et acceptée.\\[3pt]
Lisez le mot qui suit le code d'erreur avant de réessayer.
\end{piege}

\subsection{Étape 5. L'installation complète (15 min, les deux)}

Le membre B clone après avoir accepté l'invitation ; le membre A reste dans son dossier. Chacun lance
ensuite le script de mise en place.

\begin{lstlisting}[language=bash]
git clone https://github.com/votre-compte/sentiment-app.git
cd sentiment-app
powershell -ExecutionPolicy Bypass -File scripts\setup.ps1    # Windows
bash scripts/setup.sh                                         # Linux, macOS, Git Bash
\end{lstlisting}

Le script enchaîne quatre étapes affichées à l'écran : vérifier Python, créer \co{.venv}, installer
\co{requirements.txt} puis figer l'état exact dans \co{requirements.lock}, lancer le diagnostic. Ses
premières lignes méritent un commentaire :

\begin{lstlisting}[language=bash]
set -euo pipefail                      # s'arrête à la première erreur
cd "$(dirname "$0")/.."                # remonte à la racine du dépôt, d'où qu'on le lance
PY=$(command -v python3.11 || command -v python3 || command -v python)
[ -d .venv ] || "$PY" -m venv .venv    # ne recrée pas le venv s'il existe
\end{lstlisting}

\begin{center}
\begin{tabular}{@{}p{3.4cm}p{3.6cm}p{4.0cm}p{4.0cm}@{}}
\toprule
\textbf{Fichier} & \textbf{Qui l'écrit} & \textbf{Contenu} & \textbf{Dans Git ?} \\
\midrule
\co{requirements.txt} & L'enseignante & \co{pandas==2.2.*} & Oui, c'est le contrat \\
\co{requirements.lock} & \co{pip freeze}, chez vous & \co{pandas==2.2.3} & Oui, par \textbf{un seul} membre \\
\bottomrule
\end{tabular}
\end{center}

Le lock diffère légèrement d'un système à l'autre : si les deux membres le poussent, vous obtenez un
conflit à chaque fois.

\begin{controle}{T4 : le diagnostic}
\co{python scripts/check_setup.py} doit afficher \textbf{0 erreur(s)}. Un ou deux avertissements sont
normaux à ce stade : les fiches si la partie A n'est pas finie, et le fait qu'un seul membre ait
commité.
\end{controle}

Le diagnostic distingue trois niveaux. \co{OK} ne demande rien. \co{ATTENTION} ne bloque pas, mais
doit disparaître avant la fin de la séance. \co{ERREUR} bloque : venv non activé, paquet absent,
dossier manquant, pas de dépôt Git, \co{.venv/} suivi.

\subsection{Étape 6. Un commit par membre (10 min, chacun)}

\begin{lstlisting}[language=bash]
git pull                       # récupère les commits de l'autre AVANT de pousser
# ajoutez votre ligne au tableau du binôme dans README.md
git add README.md              # un seul fichier nommé, pas -A
git commit -m "README : ajout de <nom> au tableau du binome"
git push
\end{lstlisting}

Pourquoi nommer le fichier plutôt que \co{-A} : cela évite d'embarquer le \co{requirements.lock}
régénéré sur votre poste, qui provoquerait un conflit chez l'autre membre. Les messages de commit
sont notés : « README : ajout de Sara au tableau » plutôt que « update ».

En cas de conflit, après \co{git pull}, le fichier contient des marqueurs :

\begin{lstlisting}[language=bash]
<<<<<<< HEAD
| Membre A | Sara | @sara | |
=======
| Membre B | Omar | @omar | |
>>>>>>> origin/main
\end{lstlisting}

Gardez les deux lignes, supprimez les trois marqueurs, puis \co{git add README.md},
\co{git commit}, \co{git push}. Un conflit n'est pas une erreur : Git refuse simplement de choisir à
votre place.

\begin{controle}{T5 : les deux membres sont comptés}
\co{git shortlog -sn --all} affiche deux noms distincts, puis
\co{python -m pytest tests/test_git_history.py -q} donne \textbf{3 passed}.
\end{controle}

Si le test annonce \co{1 auteur(s) de commits, 2 attendus}, deux causes : un membre n'a pas encore
poussé, ou les deux postes partagent le même \co{user.name}. Dans le second cas, corrigez l'identité
puis refaites un commit ; l'ancien gardera l'ancien nom, car un commit est immuable.

\subsection{Étapes 7 et 8. Vérification et clone frais (15 min)}

\begin{lstlisting}[language=bash]
git pull
python scripts\check_setup.py      # T6 : 0 erreur(s), 0 avertissement(s)
python -m pytest -q                # T6 : 38 passed
\end{lstlisting}

Les 38 tests se répartissent en 23 d'environnement, 12 de cadrage et 3 d'historique Git. Un seul
avertissement reste toléré jusqu'au bout : celui sur la version de Python, si votre poste n'est pas
en 3.11. Poussez, puis vérifiez que l'onglet \textbf{Actions} est vert.

\begin{lstlisting}[language=bash]
bash scripts/verif_clone.sh https://github.com/votre-compte/sentiment-app.git
\end{lstlisting}

\begin{controle}{T7 : le dépôt tient debout tout seul}
\textbf{0 erreur(s)} et \textbf{38 passed} sur un clone frais, sans aucune intervention manuelle :
aucun fichier à créer, aucune commande à ajouter, aucun paquet à installer en plus.
\end{controle}

Si quelque chose manque ici alors que tout marchait dans votre dossier de travail, c'est qu'un fichier
n'a jamais été poussé. C'est exactement ce que fera l'enseignant pour noter.

% =====================================================================
\section{TP2 : entraîner, évaluer, sauvegarder le modèle}
% =====================================================================

\begin{objectif}
Produire un modèle reproductible, évalué correctement, et sauvegardé dans un fichier que l'interface
du TP3 pourra charger. Pas un bon modèle : un modèle dont vous savez pourquoi il se trompe.
\end{objectif}

La fiche de cadrage du cas 3, rédigée au TP1, devient le cahier des charges : approche machine
learning, métrique de rappel sur la classe négative, seuil d'acceptation.

\subsection{Les quatre scripts du kit}

\begin{center}
\begin{longtable}{@{}p{3.4cm}p{6.2cm}p{5.8cm}@{}}
\toprule
\textbf{Script} & \textbf{Ce qu'il fait} & \textbf{Ce qu'il produit} \\
\midrule
\endhead
\co{src/make_demo.py} & 290 avis déterministes : 80 nettement positifs, 80 nettement négatifs, 100 mitigés, 30 neutres & \co{data/demo_reviews.csv} \\
\co{src/train.py} & Découpe, entraîne un pipeline TF-IDF et régression logistique, mesure sur le jeu de test & \co{models/sentiment.joblib}, \co{models/metrics.json} \\
\co{src/evaluate.py} & Matrice de confusion, rapport de classification, trois erreurs affichées & une lecture, pas un fichier \\
\co{src/predict.py} & La fonction \co{predict()} appelée par l'interface, et plus tard par l'API & l'interface publique du modèle \\
\bottomrule
\end{longtable}
\end{center}

\subsection{Étape A. Installer le kit (15 min)}

\begin{lstlisting}[language=bash]
cd tp2-semaine2
python appliquer_tp2.py ../sentiment-app --simulation
python appliquer_tp2.py ../sentiment-app
cd ../sentiment-app && source .venv/bin/activate
python -m pytest -q
\end{lstlisting}

\begin{controle}{C1 : le point de départ}
\textbf{2 failed, 38 passed, 2 skipped}. Les 38 verts sont ceux du TP1, inchangés. Les 2 échecs
réclament \co{models/metrics.json}, qui n'existe pas encore. Les 2 ignorés attendent le modèle
lui-même.
\end{controle}

Comme au TP1, le rouge au démarrage est voulu : il décrit exactement le travail de la séance.

\subsection{Étape B. Le jeu de démonstration (15 min)}

\begin{lstlisting}[language=bash]
python src/make_demo.py
\end{lstlisting}

\begin{controle}{C2 : le jeu est là}
\textbf{290 avis, 145 positifs et 145 négatifs}, écrits dans \co{data/demo_reviews.csv}.
\end{controle}

Ouvrez le fichier et lisez trois avis mitigés avant d'aller plus loin. Ils ont tous la même forme :

\begin{lstlisting}[language=bash]
Service impeccable et livraison en avance sur le papier, sauf que ensuite
produit cassé à l'arrivée.                                        -> negative
\end{lstlisting}

Une première moitié d'un sentiment, un retournement (« mais finalement », « sauf que ensuite »), une
seconde moitié de l'autre. L'étiquette suit la \textbf{seconde} moitié, celle qui conclut. Un humain
lit cela sans hésiter. Retenez cette structure : elle explique chacune des erreurs de l'étape D.

Le jeu est déterministe, graine fixée à 42 : vous et votre binôme obtenez exactement les mêmes
290 avis, donc les mêmes scores. Même principe que le \co{random_state} du découpage.

\subsection{Étape C. Entraîner et sauvegarder (25 min)}

\begin{lstlisting}[language=bash]
python src/train.py
\end{lstlisting}

\begin{lstlisting}[language=bash]
Source : demo, 290 avis (145 négatifs)
232 avis d'entraînement, 58 de test
Vocabulaire : 316 termes
F1 macro sur le jeu de test : 0.879 (seuil 0.85)
Modèle : models/sentiment.joblib
Métriques : models/metrics.json (versionné dans Git, contrairement au modèle)
\end{lstlisting}

\begin{controle}{C3 : deux fichiers, deux destins}
\co{git status --short} doit montrer \co{models/metrics.json}, et \textbf{pas}
\co{models/sentiment.joblib}, ni \co{data/demo_reviews.csv}.
\end{controle}

La règle est la même pour les deux absents : ce qui se régénère par une commande n'entre pas dans
Git ; seule la \textbf{mesure} est versionnée. C'est ce qui permettra à l'intégration continue de
vérifier la qualité du modèle sans jamais le télécharger (section~\ref{sec:mlops}).

\subsubsection{Les trois lignes de code à lire}

\begin{lstlisting}[language=Python]
X_train, X_test, y_train, y_test = train_test_split(
    df["review"], df["label"], test_size=0.2, random_state=42, stratify=df["label"])
\end{lstlisting}

\textbf{Le découpage, avant toute vectorisation.} Si vous appreniez le vocabulaire TF-IDF sur
l'ensemble des données, les mots du jeu de test influenceraient l'entraînement. Le score monterait,
et il serait faux. C'est la fuite de données, l'erreur la plus fréquente et la plus difficile à
repérer après coup. \co{stratify} conserve la proportion des deux classes dans chaque moitié.

\begin{lstlisting}[language=Python]
Pipeline([
    ("tfidf", TfidfVectorizer(ngram_range=(1, 2), min_df=2, sublinear_tf=True)),
    ("clf", LogisticRegression(max_iter=1000, class_weight="balanced", random_state=42)),
])
\end{lstlisting}

\textbf{Le pipeline, un seul objet.} Vectoriseur et modèle voyagent ensemble : sauvegarder le
classifieur seul donnerait un objet inutilisable, privé de son vocabulaire. C'est aussi ce qui
garantit que la transformation appliquée en production est exactement celle de l'entraînement.
\co{ngram_range=(1,2)} ajoute les paires de mots aux mots isolés, \co{min_df=2} écarte les termes
vus une seule fois, \co{class_weight="balanced"} compense un éventuel déséquilibre des classes.

\textbf{La graine.} \co{random_state=42} apparaît au découpage et dans le modèle. Sans elle, deux
exécutions sur les mêmes données donnent deux scores, et aucune comparaison ne veut plus rien dire.

\subsection{Étape D. Évaluer et comprendre les erreurs (30 min)}

C'est l'étape la plus importante de la séance. Entraîner prend trois lignes ; savoir ce que vaut le
résultat prend une demi-heure.

\begin{lstlisting}[language=bash]
python src/evaluate.py
\end{lstlisting}

\begin{lstlisting}[language=bash]
              precision    recall  f1-score   support

    negative      0.844     0.931     0.885        29
    positive      0.923     0.828     0.873        29

    accuracy                          0.879        58

Matrice de confusion (lignes = vérité, colonnes = prédiction)
            negative  positive
  negative        27         2
  positive         5        24
\end{lstlisting}

\begin{controle}{C4 : la matrice de confusion}
\textbf{7 erreurs sur 58}, et les trois exemples affichés par le script sont tous des avis mitigés.
\end{controle}

La matrice se lit par ligne : sur 29 avis réellement négatifs, 27 ont été reconnus et 2 pris pour des
positifs ; sur 29 positifs, 24 reconnus et 5 pris pour des négatifs.

\begin{center}
\begin{tabular}{@{}cp{3.4cm}p{9.6cm}@{}}
\toprule
\textbf{Chiffre} & \textbf{Nom} & \textbf{Ce que ça coûte dans notre cas} \\
\midrule
27 & Vrais négatifs & Avis mécontents correctement remontés \\
2 & Faux positifs & Un client mécontent traité en retard : \textbf{le cas qui compte} \\
5 & Faux négatifs & Un avis positif remonté pour rien, sans gravité \\
24 & Vrais positifs & Rien à faire \\
\bottomrule
\end{tabular}
\end{center}

L'exactitude de 0,879 ne dit rien de tout cela. Ce sont les deux faux positifs qui pèsent, parce que
la fiche de cadrage visait le rappel sur la classe négative. \textbf{Une métrique unique cache
toujours la question qu'on se pose vraiment.}

\begin{piege}{Pourquoi ces sept erreurs}
Le modèle compte des mots positifs et des mots négatifs, puis tranche en faveur du camp le plus
fourni. Il ne voit pas le retournement, parce qu'un sac de mots \textbf{n'a pas d'ordre} : « bon au
début, mauvais ensuite » et « mauvais au début, bon ensuite » lui donnent exactement le même vecteur.
\\[4pt]
C'est la limite annoncée en cours, et vous venez de la mesurer plutôt que de la croire sur parole.
Il faudrait des représentations qui tiennent compte de l'ordre, les embeddings, pour aller plus loin.
\end{piege}

Créez \co{docs/semaine2.md} et répondez en trois phrases : quelles erreurs, pourquoi, et ce qu'il
faudrait pour les corriger. C'est noté, et c'est la partie du TP qu'un correcteur lit en premier.

\subsection{Étape E. Le vrai jeu, les tests, le push (35 min)}

\begin{lstlisting}[language=bash]
python src/train.py --allocine
\end{lstlisting}

Le script télécharge 20 000 avis de films en français depuis une URL publique (jeu Allociné, 7,6 Mo,
quelques secondes) et réentraîne le même pipeline, sans qu'une ligne de code change. Si le réseau est
coupé ou filtré, il bascule seul sur le jeu de démonstration et le signale.

\begin{center}
\begin{tabular}{@{}p{5.4cm}rr@{}}
\toprule
 & \textbf{Démonstration} & \textbf{Allociné} \\
\midrule
Avis & 290 & 20 000 \\
Vocabulaire & 316 termes & 149 529 termes \\
F1 macro & 0,879 & \textbf{0,920} \\
Rappel sur la classe négative & 0,931 & 0,911 \\
\bottomrule
\end{tabular}
\end{center}

Le modèle n'a pas changé : ce sont les données. Un vocabulaire 470 fois plus riche, et surtout une
mesure enfin crédible, calculée sur 4 000 avis et non sur 58. Avec 58 avis de test, une erreur de
plus déplace le score de deux points : on ne conclut rien d'un tel échantillon.

\begin{controle}{C5 : tout est vert}
\co{python -m pytest -q} affiche \textbf{42 passed}, soit les 38 du TP1 plus les 4 du modèle.
Poussez, puis vérifiez que l'onglet Actions est vert.
\end{controle}

% =====================================================================
\section{TP3 : l'interface de démonstration}
% =====================================================================

\begin{objectif}
Rendre le modèle utilisable par quelqu'un qui n'ouvrira jamais un terminal, sans réécrire une ligne
de logique de prédiction.
\end{objectif}

\subsection{Streamlit en trois idées}

\textbf{Un script Python devient une page web.} Pas de HTML, pas de JavaScript, pas de serveur à
configurer : chaque appel à \co{st.title}, \co{st.text_area} ou \co{st.dataframe} ajoute un élément
à la page.

\textbf{Chaque interaction réexécute le script de haut en bas.} C'est le modèle de fonctionnement, et
il surprend toujours. D'où le décorateur \co{@st.cache_resource} sur le chargement du modèle : sans
lui, le fichier \co{joblib} serait relu à chaque clic.

\textbf{C'est une démonstration, pas l'application.} Elle sert à montrer le modèle, à récolter des
retours, à faire apparaître les cas où il se trompe. C'est l'API, plus tard, qui sera appelée par
d'autres programmes.

\subsection{Pourquoi deux fichiers}

\begin{lstlisting}[language=Python]
# ui/helpers.py : la logique, testable sans Streamlit
def classer(texte, predicteur):
    resultat = predicteur(texte)
    resultat["incertain"] = resultat["score"] < SEUIL_INCERTITUDE
    resultat["verdict"] = "à vérifier" if resultat["incertain"] else resultat["label"]
    return resultat
\end{lstlisting}

\begin{lstlisting}[language=Python]
# ui/app.py : l'affichage, et rien d'autre
r = classer(texte, predire)
if r["incertain"]:
    st.warning("À vérifier par un agent")
elif r["label"] == "negative":
    st.error("Avis négatif : à traiter en priorité")
else:
    st.success("Avis positif")
\end{lstlisting}

Les trois tests du TP3 importent \co{helpers.py} \textbf{sans lancer Streamlit}, et c'est ce que
vérifie \co{test_helpers_importables}. Une interface non testable est une interface qu'on casse sans
s'en apercevoir.

\subsection{Étape A. Installer et lancer (15 min)}

\begin{lstlisting}[language=bash]
cd tp3-streamlit
python appliquer_tp3.py ../sentiment-app
cd ../sentiment-app && source .venv/bin/activate
streamlit run ui/app.py
\end{lstlisting}

\begin{controle}{S1 : la page s'ouvre}
Le navigateur s'ouvre sur \co{http://localhost:8501} et affiche trois onglets. Si le modèle manque,
la page le dit et s'arrête : lancez \co{python src/train.py --allocine}, puis rechargez.
\end{controle}

Ce terminal reste occupé tant que l'interface tourne : ouvrez-en un second pour la suite, et arrêtez
avec \co{Ctrl+C}.

\subsection{Étape B. L'onglet « un avis » (30 min)}

Saisissez trois avis : un clairement positif, un clairement négatif, puis un mitigé du type « service
impeccable au début, mais finalement produit cassé ».

\begin{controle}{S2 : le modèle dit quand il hésite}
L'avis mitigé doit afficher \textbf{« à vérifier »} plutôt qu'un verdict tranché, dès que la
confiance tombe sous 60 \%.
\end{controle}

Le seuil est dans \co{ui/helpers.py}, constante \co{SEUIL_INCERTITUDE}. Faites-le varier et notez ce
qui change :

\begin{center}
\begin{tabular}{@{}cp{4.4cm}p{8.6cm}@{}}
\toprule
\textbf{Seuil} & \textbf{Avis envoyés à un agent} & \textbf{Risque} \\
\midrule
0,50 & peu & Le modèle tranche seul des cas qu'il ne maîtrise pas \\
0,60 & modéré & Compromis retenu par défaut \\
0,75 & beaucoup & La file humaine se remplit, l'automatisation perd son intérêt \\
\bottomrule
\end{tabular}
\end{center}

Ce curseur ne se règle pas dans le code : il se décide avec l'équipe métier, et il découle directement
de la fiche du cas 3, qui disait qu'un avis mal classé se requalifie d'un clic.

\subsection{Étape C. L'onglet « un fichier CSV » (40 min)}

Téléversez \co{data/sample_reviews.csv}, puis \co{data/demo_reviews.csv} produit au TP2.

\begin{controle}{S3 : le résumé et le tableau}
Trois métriques en haut (avis analysés, négatifs détectés avec leur pourcentage, cas à vérifier), un
graphique de répartition, un tableau trié négatifs en tête, et un bouton de téléchargement.
\end{controle}

Sur \co{demo_reviews.csv}, comparez la part de négatifs annoncée par l'interface avec le vrai
étiquetage du fichier : l'écart correspond exactement aux erreurs mesurées au TP2.

\subsection{Étape D. Tests, documentation, push (35 min)}

\begin{lstlisting}[language=bash]
python -m pytest tests/test_ui.py -q      # 3 passed
python -m pytest -q                       # 45 passed
git add -A
git commit -m "Semaine 3 : interface de demonstration Streamlit"
git push
\end{lstlisting}

\begin{controle}{S4 : tout est vert}
\textbf{45 passed}, soit 38 du TP1, 4 du TP2 et 3 du TP3, puis la CI verte.
\end{controle}

Écrivez \co{docs/semaine3.md}, trois phrases : quel seuil vous retenez, combien d'avis cela envoie à
un agent sur le jeu de démonstration, et pourquoi ce compromis vous paraît juste.

\subsection{Sur Google Colab}

Streamlit ne s'affiche pas directement dans Colab : le serveur tourne sur la machine virtuelle, que
votre navigateur ne peut pas joindre. Le notebook \co{TP3-colab-streamlit.ipynb} passe donc par un
tunnel.

\begin{lstlisting}[language=bash]
!curl -s https://loca.lt/mytunnelpassword     # le mot de passe du tunnel : l'IP publique
!streamlit run ui/app.py --server.headless true --server.port 8501 &>/content/log &
!npx --yes localtunnel --port 8501            # affiche une adresse loca.lt à ouvrir
\end{lstlisting}

\begin{piege}{Trois différences avec un poste local}
\textbf{Le tunnel} donne une adresse \textbf{publique et temporaire} : aucune donnée réelle ne doit y
transiter.\\[3pt]
\textbf{Pas d'environnement virtuel} : Colab installe dans son Python système, donc
\co{test_virtualenv_active} est écarté et les comptes deviennent 37, 41 et 44 au lieu de 38, 42
et 45.\\[3pt]
\textbf{Rien ne survit à la session} : \co{/content} disparaît à la fermeture. Poussez sur GitHub
avant de partir. Le modèle, lui, n'a pas à être conservé.
\end{piege}

% =====================================================================
\section{Les 45 tests}
% =====================================================================

\begin{objectif}
Comprendre ce que la machine vérifie, pour corriger seul au lieu d'attendre.
\end{objectif}

Le dépôt se vérifie de deux façons, qui contrôlent les mêmes choses. \co{scripts/check_setup.py}
parle aux humains, avec des lignes OK, ATTENTION, ERREUR. \co{pytest} parle aux machines, et c'est
lui que la CI et la notation utilisent.

\subsection{Répartition}

\begin{center}
\begin{tabular}{@{}p{4.4cm}cp{9.0cm}@{}}
\toprule
\textbf{Fichier} & \textbf{Tests} & \textbf{Ce qu'il vérifie} \\
\midrule
\co{test_environment.py} & 23 & Version Python, venv actif, trois paquets, sept dossiers, quatre fichiers, versions figées, cinq motifs du \co{.gitignore}, rien de secret commité \\
\co{test_cadrage.py} & 12 & Trois fiches, huit rubriques, rubriques remplies, une seule case cochée, deux critères cités \\
\co{test_git_history.py} & 3 & Au moins un commit, deux auteurs distincts, \co{.venv/} absent de tout l'historique \\
\co{test_model.py} & 4 & Métriques complètes, F1 au-dessus du seuil, prédiction correcte, format de sortie \\
\co{test_ui.py} & 3 & Helpers importables sans Streamlit, verdict ajouté, traitement d'un lot \\
\bottomrule
\end{tabular}
\end{center}

\subsection{Ce qui est ignoré, et pourquoi}

Certains tests s'effacent au lieu d'échouer quand leur condition d'exécution n'est pas réunie. C'est
voulu, et cela explique les \co{skipped} des contrôles T1 et C1.

\begin{center}
\begin{tabular}{@{}p{5.6cm}p{9.8cm}@{}}
\toprule
\textbf{Tests} & \textbf{Ignorés quand} \\
\midrule
Les 3 de \co{test_git_history.py} & Git manque, \co{.git/} n'existe pas, ou aucun commit \\
\co{test_no_real_env_file_committed} & Le dossier n'est pas un dépôt Git \\
Les 2 tests de prédiction du modèle & \co{models/sentiment.joblib} absent, ce qui est le cas en CI \\
Les 2 tests d'interface utilisant le modèle & Idem \\
\bottomrule
\end{tabular}
\end{center}

\subsection{Les paramètres réglables}

\begin{center}
\begin{tabular}{@{}p{4.6cm}p{5.4cm}p{2.0cm}p{3.0cm}@{}}
\toprule
\textbf{Paramètre} & \textbf{Fichier} & \textbf{Valeur} & \textbf{Rôle} \\
\midrule
\co{MIN_AUTHORS} & \co{check_setup.py}, \co{test_git_history.py} & 2 & Auteurs exigés \\
\co{REQUIRED_DIRS} & \co{check_setup.py}, \co{test_environment.py} & 7 & La structure attendue \\
\co{REQUIRED_SECTIONS} & \co{test_cadrage.py} & 8 & Rubriques des fiches \\
\co{CRITERIA} & \co{test_cadrage.py} & (a) à (d) & Lettres acceptées \\
\co{SEUIL_F1} & \co{train.py}, \co{test_model.py} & 0,85 & Seuil d'acceptation \\
\co{SEUIL_INCERTITUDE} & \co{ui/helpers.py} & 0,60 & Bascule vers « à vérifier » \\
\bottomrule
\end{tabular}
\end{center}

Ces constantes sont en haut de chaque fichier, lisibles sans chercher. \textbf{Ne les modifiez pas
pour faire passer un test} : c'est exactement ce que la notation regarde. Le seul qu'on vous demande
de faire varier est \co{SEUIL_INCERTITUDE}, et c'est l'objet du contrôle S2.

% =====================================================================
\section{Intégration continue et MLOps}
\label{sec:mlops}
% =====================================================================

\subsection{Trois artefacts, pas un}

Un logiciel classique versionne une chose : le code. Une application intelligente en versionne trois,
et ce tableau est la règle qui décide, pour chacune, où elle vit.

\begin{center}
\begin{tabular}{@{}p{5.0cm}p{2.6cm}p{8.0cm}@{}}
\toprule
\textbf{Artefact} & \textbf{Dans Git ?} & \textbf{Pourquoi} \\
\midrule
Le code (\co{src/}, \co{ui/}, \co{tests/}) & Oui & C'est le travail \\
L'environnement (\co{.venv/}) & Non & Reconstructible : \co{pip install -r requirements.txt} \\
Les données (\co{demo_reviews.csv}) & Non & Reconstructible : \co{python src/make_demo.py} \\
Le modèle (\co{sentiment.joblib}) & Non & Reconstructible : \co{python src/train.py} \\
\textbf{La mesure} (\co{metrics.json}) & \textbf{Oui} & Elle ne se reconstruit pas sans réentraîner \\
\bottomrule
\end{tabular}
\end{center}

\subsection{Ce que fait l'intégration continue}

\begin{lstlisting}[language=bash]
on: [push, pull_request]
jobs:
  tp1:
    runs-on: ubuntu-latest
    steps:
      - uses: actions/checkout@v4
        with: { fetch-depth: 0 }
      - uses: actions/setup-python@v5
        with: { python-version: "3.11" }
      - run: python -m venv .venv && . .venv/bin/activate && pip install pandas scikit-learn pytest
      - run: . .venv/bin/activate && python scripts/check_setup.py
      - run: . .venv/bin/activate && pytest -q
\end{lstlisting}

\begin{center}
\begin{tabular}{@{}p{4.6cm}p{11.0cm}@{}}
\toprule
\textbf{Paramètre} & \textbf{Rôle} \\
\midrule
\co{on: [push, pull_request]} & Les tests tournent à chaque push et à chaque pull request \\
\co{fetch-depth: 0} & Récupère \textbf{tout l'historique}. Par défaut Actions ne clone qu'un commit, et \co{git shortlog} ne verrait qu'un auteur : le test échouerait toujours \\
\co{python-version: "3.11"} & La version de référence du cours, quelle que soit la vôtre \\
Création du venv dans chaque \co{run} & Sinon \co{test_virtualenv_active} échouerait \\
Trois paquets seulement & La CI dure une minute ; l'installation complète est vérifiée au clone frais \\
\bottomrule
\end{tabular}
\end{center}

Une CI verte veut dire que votre dépôt fonctionne \textbf{ailleurs que sur vos deux portables}. Si
elle est rouge alors que tout passe chez vous, ouvrez le journal : c'est le même pytest, et le nom du
test en échec dit quoi corriger.

\subsection{Le geste MLOps de la séquence}

\begin{piege}{La CI n'a jamais votre modèle}
\co{models/sentiment.joblib} n'est pas dans le dépôt. Pourtant, l'intégration continue vérifie qu'un
modèle acceptable a été produit.
\end{piege}

\begin{center}
\begin{tabular}{@{}p{7.0cm}p{8.6cm}@{}}
\toprule
\textbf{Les 4 tests du modèle} & \textbf{En intégration continue} \\
\midrule
Deux lisent \co{models/metrics.json} & Ils passent : le fichier est versionné \\
Deux appellent \co{predict()} & Ils sont ignorés : le modèle est absent \\
\bottomrule
\end{tabular}
\end{center}

On transporte la \textbf{preuve} de la qualité, pas l'artefact lourd. C'est le principe que MLflow et
DVC industrialiseront plus tard dans le semestre, et la porte de validation qui déclenchera le
réentraînement automatique quand la dérive sera au programme.

\subsection{DevOps, MLOps : où passe la frontière}

\begin{center}
\begin{tabular}{@{}p{7.4cm}p{8.2cm}@{}}
\toprule
\textbf{DevOps, dans ces trois TP} & \textbf{MLOps, ce qui s'y ajoute} \\
\midrule
Environnement reproductible & Trois artefacts au lieu d'un \\
Tests automatisés & Une validation qui compare une performance à un seuil, au lieu de passer ou d'échouer \\
Versionnage à plusieurs & Des données et un modèle à versionner, trop lourds pour Git \\
Intégration continue & Une dégradation possible \textbf{sans qu'on ait rien modifié} : la dérive \\
\bottomrule
\end{tabular}
\end{center}

Un test unitaire passe ou échoue. Un modèle a un F1 à comparer à un seuil, sur un jeu de test figé :
c'est pour cela que \co{metrics.json} existe, et que le seuil y est inscrit à côté de la mesure.

% =====================================================================
\section{Le compte rendu global}
% =====================================================================

À la fin des trois séances, chaque binôme rédige un compte rendu unique, de quatre à six pages,
déposé dans son dépôt et figé par une étiquette Git.

\begin{lstlisting}[language=bash]
cp docs/compte-rendu-TEMPLATE.md docs/compte-rendu.md
# rédaction...
git add docs/compte-rendu.md
git commit -m "Compte rendu global des TP1 a TP3"
git push
git tag -a compte-rendu -m "Version remise le <date>"
git push origin compte-rendu
\end{lstlisting}

L'étiquette fige la version remise : les commits postérieurs ne sont pas pris en compte. Format
Markdown, dans le dépôt, pas de PDF ni de lien externe : le compte rendu vit avec le code qu'il
décrit.

\begin{center}
\begin{longtable}{@{}p{5.6cm}cp{8.4cm}@{}}
\toprule
\textbf{Section} & \textbf{Pts} & \textbf{Ce qui est regardé} \\
\midrule
\endhead
1. Le besoin et son cadrage & 3 & Les lettres (a) à (d) comme raisonnement, pas comme formalité \\
2. L'architecture du dépôt & 3 & Vous savez dire pourquoi \co{metrics.json} est dans Git et pas le modèle \\
3. Le modèle & 4 & Vos F1 correspondent à votre \co{metrics.json} ; le découpage avant vectorisation est expliqué \\
4. L'analyse des erreurs & 4 & Les avis mitigés et l'absence d'ordre, avec un exemple tiré de \textbf{vos} données \\
5. Le seuil d'incertitude & 3 & Combien d'avis partent vers un agent, et pourquoi ce compromis \\
6 à 8. Échecs, recul, répartition & 3 & Un compte rendu sans aucun échec n'est pas crédible \\
\bottomrule
\end{longtable}
\end{center}

Trois règles : citez \textbf{vos} chiffres et non ceux du corrigé, avec une phrase sur l'écart le cas
échéant ; un seul compte rendu par binôme, avec une section sincère sur la répartition du travail ;
et si vous utilisez un assistant d'IA pour rédiger, dites-le en une phrase, en précisant ce que vous
avez vérifié vous-même. C'est autorisé, le cacher ne l'est pas.

% =====================================================================
\section{Récapitulatif}
% =====================================================================

\subsection{Les dix-sept contrôles}

\begin{center}
\begin{longtable}{@{}p{1.0cm}p{6.8cm}p{7.4cm}@{}}
\toprule
 & \textbf{Commande} & \textbf{Résultat attendu} \\
\midrule
\endhead
T0 & \co{ls} et \co{ls -a} & 7 dossiers, \co{.gitignore} présent \\
T1 & \co{pytest -q} & 3 failed, 31 passed, 4 skipped \\
T2 & \co{pytest tests/test_cadrage.py -q} & 12 passed \\
T3 & \co{git status} avant le commit & aucune ligne \co{.venv/} \\
T4 & \co{python scripts/check_setup.py} & 0 erreur(s) \\
T5 & \co{git shortlog -sn --all} puis les tests Git & deux noms, 3 passed \\
T6 & diagnostic et \co{pytest -q} & 0 erreur, 0 avertissement, 38 passed \\
T7 & la suite sur un clone frais & 38 passed sans intervention \\
C1 & \co{pytest -q} après le kit TP2 & 2 failed, 38 passed, 2 skipped \\
C2 & \co{python src/make_demo.py} & 290 avis, 145 de chaque classe \\
C3 & \co{python src/train.py} & F1 macro 0,879 ; seul \co{metrics.json} en attente \\
C4 & \co{python src/evaluate.py} & 7 erreurs sur 58, toutes sur des avis mitigés \\
C5 & \co{train.py --allocine} puis \co{pytest -q} & F1 0,920 ; 42 passed ; CI verte \\
S1 & \co{streamlit run ui/app.py} & la page s'ouvre, trois onglets \\
S2 & un avis mitigé & verdict « à vérifier » sous 60 \% \\
S3 & un CSV dans l'onglet 2 & métriques, graphique, tableau trié, export \\
S4 & \co{pytest -q} & 45 passed, CI verte \\
\bottomrule
\end{longtable}
\end{center}

\subsection{Les erreurs les plus coûteuses}

\begin{center}
\begin{longtable}{@{}p{6.0cm}p{4.4cm}p{4.8cm}@{}}
\toprule
\textbf{Message} & \textbf{Cause} & \textbf{Correction} \\
\midrule
\endhead
\co{pytest n'est pas reconnu} & Venv non activé & Activer, ou \co{python -m pytest} \\
\co{Activate.ps1 cannot be loaded} & Politique PowerShell & \co{Set-ExecutionPolicy -Scope CurrentUser RemoteSigned} \\
\co{ModuleNotFoundError} & Venv non activé ou installation partielle & \co{where python} doit contenir \co{.venv} \\
\co{No matching distribution found} & Python 3.12 ou 3.13 & Installer 3.11, recréer le venv \\
\co{error: 400} sur l'URL & Chevrons restés dans l'adresse & \co{git remote set-url origin <vraie URL>} \\
\co{Password authentication is not supported} & Authentification & Jeton d'accès personnel \\
\co{403 Forbidden} & Autorisation & Accepter l'invitation de collaborateur \\
\co{rejected, non-fast-forward} & L'autre membre a poussé & \co{git pull}, résoudre, \co{git push} \\
\co{1 auteur(s) de commits, 2 attendus} & Même \co{user.name} sur les deux postes & Identité distincte, puis un vrai commit \\
\co{.venv/ a été commité} & \co{.gitignore} absent au premier \co{add} & Refaire le dépôt \\
\co{models/metrics.json absent} & Étape C du TP2 sautée & \co{python src/train.py} \\
\co{models/sentiment.joblib absent} & Modèle non entraîné sur ce poste & \co{python src/train.py --allocine} \\
\co{Port 8501 is already in use} & Une interface tourne déjà & \co{Ctrl+C}, ou \co{--server.port 8502} \\
\bottomrule
\end{longtable}
\end{center}

\subsection{La structure finale du dépôt}

\begin{lstlisting}[language=bash]
sentiment-app/
├── README.md TP1-ENONCE.md MARCHE-A-SUIVRE.md TESTS.md AIDE-GIT.md DEPANNAGE.md
├── requirements.txt requirements.lock .gitignore .env.example pytest.ini Makefile
├── .github/workflows/ci.yml        intégration continue
├── docs/
│   ├── cadrage/cas1, cas2, cas3    TP1, remplies (livrable noté)
│   ├── semaine2.md                 TP2, analyse des erreurs
│   ├── semaine3.md                 TP3, seuil d'incertitude
│   └── compte-rendu.md             le rendu final, figé par une étiquette Git
├── src/      make_demo, train, evaluate, predict          (TP2)
├── ui/       app.py, helpers.py                           (TP3)
├── tests/    environment (23), cadrage (12), git (3), model (4), ui (3)
├── scripts/  check_setup, setup, init_depot, verif_clone, new_cadrage
├── models/   .gitkeep + metrics.json    le modèle n'est JAMAIS commité
└── data/     sample_reviews.csv         demo_reviews.csv est régénérable
\end{lstlisting}

\textbf{45 tests} à l'arrivée. Trois artefacts restent hors de Git parce qu'ils se reconstruisent en
une commande : le venv, le jeu de démonstration, le modèle. Une seule mesure est versionnée. C'est le
principe qui portera tout le reste du semestre, et le dépôt d'aujourd'hui est celui de la dernière
semaine.

\end{document}
